\section{Consensus: from \PoW to \PoN, node tiering, benchmarking, and eligibility}
\label{sec:consensus}

For roughly the first two million blocks of its history, \Flux's base-layer daemon, \fluxd, produced
blocks the way most Zcash-lineage chains do: miners raced to solve an Equihash puzzle, and the chain
followed whichever fork carried the most accumulated hashing work. Since height 2{,}020{,}000 it does
not. Block production is instead awarded, once every 30 seconds, to whichever registered \FluxNode
wins a per-slot lottery keyed by its collateral outpoint. \Flux calls this mechanism Proof-of-Node
(\PoN). The name suggests that operating a node is what earns the right to produce a block. The
mechanics below show something narrower: the right is earned by a node that has posted the correct
collateral amount on-chain and that has, at some point, produced one signed message from its own
locally-run benchmarking process. Whether that local, same-host message is a hard-to-fake attestation
of real hardware or a token any operator's own machine can produce for itself is the question this
section spends most of its length answering precisely, because the answer is not the one the name
implies. The remainder of this section is formal.

\subsection{The network-upgrade mechanism}
\label{sec:consensus-upgrades}

Network upgrades are indexed by \texttt{Consensus::\allowbreak{}Upgrade\allowbreak{}Index}
\shipped \srccode{fluxd/\allowbreak{}src/\allowbreak{}consensus/\allowbreak{}params.h}{24--40}. Upgrade state is not a persisted flag; it is
a pure function of block height, recomputed on every call by \texttt{Network\allowbreak{}Upgrade\allowbreak{}State}: an upgrade
is \texttt{UPGRADE\_DISABLED} if its activation height equals the sentinel
\texttt{NO\_ACTIVATION\_HEIGHT} (\,$-1$\,), \texttt{UPGRADE\_PENDING} below its activation height, and
\texttt{UPGRADE\_ACTIVE} at or above it \shipped \srccode{fluxd/\allowbreak{}src/\allowbreak{}consensus/\allowbreak{}upgrades.cpp}{72--97}.
For \PoN specifically this reduces to a single inequality:
\texttt{Is\allowbreak{}PONActive(n\allowbreak{}Height)} $\Leftrightarrow$ \texttt{nHeight} $\ge$ \texttt{Get\allowbreak{}PONActivation\allowbreak{}Height()}
\shipped \srccode{fluxd/\allowbreak{}src/\allowbreak{}pon/\allowbreak{}pon-fork.cpp}{13--20,45--47}. There is no signalling, no
miner- or node-voting window, and no grace period at any upgrade boundary, \PoN included: activation
is a hard height cutoff.

\begin{table}[H]
\centering
\caption{Mainnet network-upgrade heights, \texttt{CMain\allowbreak{}Params::\allowbreak{}CMain\allowbreak{}Params()}. Hashes are truncated
to 16 hex characters for display; full values are at the cited lines.}
\label{tab:pon-upgrades}
\begin{tabular}{lrll}
\toprule
Upgrade & Height & Activation block hash (truncated) & Provenance \\
\midrule
\texttt{UPGRADE\_LWMA}      & 125{,}000     & --- & \srccode{fluxd/\allowbreak{}src/\allowbreak{}chainparams.cpp}{116--117} \\
\texttt{UPGRADE\_EQUI144\_5}& 125{,}100     & --- & \srccode{fluxd/\allowbreak{}src/\allowbreak{}chainparams.cpp}{119--120} \\
\texttt{UPGRADE\_ACADIA}    & 250{,}000     & \texttt{0000001d65fa78f2\ldots} & \srccode{fluxd/\allowbreak{}src/\allowbreak{}chainparams.cpp}{122--125} \\
\texttt{UPGRADE\_KAMIOOKA}  & 372{,}500     & \texttt{00000052e2ac144c\ldots} & \srccode{fluxd/\allowbreak{}src/\allowbreak{}chainparams.cpp}{127--130} \\
\texttt{UPGRADE\_KAMATA}    & 558{,}000     & \texttt{000000a33d38f37f\ldots} & \srccode{fluxd/\allowbreak{}src/\allowbreak{}chainparams.cpp}{132--135} \\
\texttt{UPGRADE\_FLUX}      & 835{,}554     & \texttt{000000ce99aa6765\ldots} & \srccode{fluxd/\allowbreak{}src/\allowbreak{}chainparams.cpp}{137--140} \\
\texttt{UPGRADE\_HALVING}   & 1{,}076{,}532 & \texttt{000000111f8643ce\ldots} & \srccode{fluxd/\allowbreak{}src/\allowbreak{}chainparams.cpp}{142--145} \\
\texttt{UPGRADE\_P2SHNODES} & 1{,}549{,}500 & \texttt{00000009f9178347\ldots} & \srccode{fluxd/\allowbreak{}src/\allowbreak{}chainparams.cpp}{147--150} \\
\texttt{UPGRADE\_PON}       & \textbf{2{,}020{,}000} & \emph{unset} & \srccode{fluxd/\allowbreak{}src/\allowbreak{}chainparams.cpp}{152--153} \\
\bottomrule
\end{tabular}
\end{table}

\texttt{BASE\_SPROUT} (always active from genesis) and \texttt{UPGRADE\_TESTDUMMY} (never activated on
mainnet) are omitted from Table~\ref{tab:pon-upgrades} because neither carries a meaningful mainnet
height \shipped \srccode{fluxd/\allowbreak{}src/\allowbreak{}chainparams.cpp}{108--114}.

Height 2{,}020{,}000 for \texttt{UPGRADE\_PON} is a code fact \shipped \srccode{fluxd/\allowbreak{}src/\allowbreak{}chainparams.cpp}{152--153}.
The calendar date sometimes attached to this activation is not: \texttt{UPGRADE\_PON} is the only
upgrade in Table~\ref{tab:pon-upgrades} with no \texttt{hash\allowbreak{}Activation\allowbreak{}Block} set, and no timestamp
literal tying height 2{,}020{,}000 to a wall-clock date appears anywhere in the source
\shipped \srccode{fluxd/\allowbreak{}src/\allowbreak{}chainparams.cpp}{152--153}. The nearest bracketing evidence is a
checkpoint at height 2{,}022{,}000 carrying UNIX timestamp \texttt{1761701944}, which is 2025-10-29
\shipped \srccode{fluxd/\allowbreak{}src/\allowbreak{}chainparams.cpp}{281,284}; two thousand blocks at 30 s spacing is
approximately 16.7 hours, so this brackets but does not establish an activation date.
\srcmeasured{block \num{2020000} carries timestamp \num{1761415235}, i.e.\ 2025-10-25T18:00:35Z, retrieved from \texttt{api.\allowbreak{}runonflux.\allowbreak{}io/\allowbreak{}daemon/\allowbreak{}getblock}; the date is therefore a chain fact rather than an inference from a nearby height}{2026-09-04}

\subsection{What changes at \PoN activation}
\label{sec:consensus-changes}

\begin{itemize}
\item \textbf{Subsidy.} \texttt{Get\allowbreak{}Block\allowbreak{}Subsidy} branches first on \texttt{Is\allowbreak{}PONActive(n\allowbreak{}Height)},
bypassing the legacy slow-start/halving formula entirely for \PoN heights
\shipped \srccode{fluxd/\allowbreak{}src/\allowbreak{}main.cpp}{2796--2853}. For $\hgt \ge \hpon$,
\[
\subsidypon(\hgt) \;=\; 14\flux \cdot \redfac^{\,n(\hgt)}, \qquad
n(\hgt) \;=\; \min\!\Big(\floorfrac{\hgt - \hpon}{\redint},\ \redmax\Big),
\]
with $\redfac = 9/10$ applied once per elapsed reduction interval, $\redint = 1{,}051{,}200$ blocks
mainnet, and a cap of $\redmax = 20$ reductions
\shipped \srccode{fluxd/\allowbreak{}src/\allowbreak{}main.cpp}{2812}, \srccode{fluxd/\allowbreak{}src/\allowbreak{}chainparams.cpp}{156--158}. The full
emission derivation, including the tail after $\redmax$ is reached, is given in \S5.

\item \textbf{Difficulty retarget.} \texttt{Get\allowbreak{}Next\allowbreak{}Work\allowbreak{}Required\allowbreak{}By\allowbreak{}Fork} dispatches to
\texttt{Get\allowbreak{}Next\allowbreak{}PONWork\allowbreak{}Required} for \PoN heights and to the legacy staged
Digishield/LWMA/LWMA3 retarget otherwise
\shipped \srccode{fluxd/\allowbreak{}src/\allowbreak{}pon/\allowbreak{}pon-fork.cpp}{22--42}, \srccode{fluxd/\allowbreak{}src/\allowbreak{}pow.cpp}{21--97}.

\item \textbf{Block header format.} Block version $\ge$ \texttt{PON\_VERSION} $= 100$
\shipped \srccode{fluxd/\allowbreak{}src/\allowbreak{}primitives/\allowbreak{}block.h}{31} adds two consensus fields:
\texttt{nodes\allowbreak{}Collateral} (the \texttt{COutPoint} of the winning \FluxNode's collateral) and
\texttt{vchBlockSig} (an ECDSA signature over the block hash)
\shipped \srccode{fluxd/\allowbreak{}src/\allowbreak{}primitives/\allowbreak{}block.h}{44--45,64--96}. \texttt{IsPON()}\slash\texttt{IsPOW()} are
simply \texttt{nVersion} $\ge$/$<$ 100 \shipped \srccode{fluxd/\allowbreak{}src/\allowbreak{}primitives/\allowbreak{}block.h}{95--96}.

\item \textbf{Block spacing.} Legacy target spacing $\spacingpow = 120\,\mathrm{s}$
\shipped \srccode{fluxd/\allowbreak{}src/\allowbreak{}chainparams.cpp}{101} drops to $\spacingpon = 30\,\mathrm{s}$
\shipped \srccode{fluxd/\allowbreak{}src/\allowbreak{}chainparams.cpp}{105}.

\item \textbf{Header validation.} \texttt{Check\allowbreak{}Block\allowbreak{}Header} replaces the Equihash-solution check and
\texttt{Check\allowbreak{}Proof\allowbreak{}Of\allowbreak{}Work} with \texttt{Check\allowbreak{}Proof\allowbreak{}Of\allowbreak{}Node} plus a signature-size cap
\shipped \srccode{fluxd/\allowbreak{}src/\allowbreak{}main.cpp}{5330--5378}. \texttt{Contextual\allowbreak{}Check\allowbreak{}Block\allowbreak{}Header} additionally
requires \texttt{block.IsPON()} to equal \texttt{Is\allowbreak{}PONActive(n\allowbreak{}Height)} exactly, rejecting a mismatch
as \texttt{"not-\allowbreak{}pon-\allowbreak{}block"} or \texttt{"premature-\allowbreak{}pon-\allowbreak{}block"}, and requires a \PoN block's timestamp to
be strictly greater than its parent's timestamp rather than median-time-past-based as the legacy path
is \shipped \srccode{fluxd/\allowbreak{}src/\allowbreak{}main.cpp}{5516--5545,5564--5572}.

\item \textbf{Future-time tolerance.} Accepted clock skew for a new block's timestamp drops from
2 hours (pre-LWMA) / 360 s (post-LWMA \PoW) to 300 s for \PoN blocks
\shipped \srccode{fluxd/\allowbreak{}src/\allowbreak{}main.cpp}{5330--5378}.

\item \textbf{Coinbase dev-fund rule.} Once \PoN is active, the coinbase must pay a minimum amount to a
hardcoded address or the transaction is rejected. That minimum is a remainder:
\[
\devfund(\hgt) \;=\; \subsidypon(\hgt) \;-\; \sum_{x \in \tiers} \tbase{x} \cdot \frac{\subsidypon(\hgt)}{14\flux},
\]
with $\tbase{\tC} = 1\flux$, $\tbase{\tN} = 3.5\flux$, $\tbase{\tS} = 9\flux$
\shipped \srccode{fluxd/\allowbreak{}src/\allowbreak{}main.cpp}{2892--2895}, which sum to $13.5\flux$ of the $14\flux$ initial
subsidy, leaving $\devfund(\hpon) = 0.5\flux$ at the first post-activation block, decaying at the same
rate as $\subsidypon$ thereafter \shipped \srccode{fluxd/\allowbreak{}src/\allowbreak{}main.cpp}{2877--2884}. This is enforced
twice: loosely (a $\ge$ check) in \texttt{Contextual\allowbreak{}Check\allowbreak{}Transaction}, which rejects a coinbase missing
it as \texttt{"tx-\allowbreak{}coinbase-\allowbreak{}missing-\allowbreak{}dev-\allowbreak{}funding"} \shipped \srccode{fluxd/\allowbreak{}src/\allowbreak{}main.cpp}{1231--1249},
and again --- also as a $\ge$ check, but against the remainder left after the exact tier payouts --- in \texttt{Fluxnode\allowbreak{}Cache::\allowbreak{}Check\allowbreak{}Fluxnode\allowbreak{}Payout}
\shipped \srccode{fluxd/\allowbreak{}src/\allowbreak{}fluxnode/\allowbreak{}fluxnode.cpp}{822--889}. The recipient is the hardcoded mainnet
address \texttt{t3h\allowbreak{}Pu1YDe\allowbreak{}GUCp8m7BQCnn\allowbreak{}NUm\allowbreak{}RMJBa5Rady\allowbreak{}A} \shipped \srccode{fluxd/\allowbreak{}src/\allowbreak{}chainparams.cpp}{306}.
\end{itemize}

\subsection{\PoN sortition, formally}
\label{sec:consensus-sortition}

\begin{definition}[\PoN slot and eligibility]
\label{def:pon-eligibility}
Let $t_0$ be network genesis time. The \PoN slot at wall-clock time $t$ is
\[
\tslot(t) \;\defeq\; \floorfrac{t - t_0}{\spacingpon}, \qquad \spacingpon = 30\,\mathrm{s}
\]
\shipped \srccode{fluxd/\allowbreak{}src/\allowbreak{}pon/\allowbreak{}pon.cpp}{88--93}. For a confirmed \FluxNode with collateral outpoint
$c$, define
\[
H(c,\tslot) \;\defeq\; \mathrm{Hash}\big(c \,\|\, h_{\mathrm{prev}} \,\|\, \tslot\big),
\]
where $h_{\mathrm{prev}}$ is the hash of the current chain tip
\shipped \srccode{fluxd/\allowbreak{}src/\allowbreak{}pon/\allowbreak{}pon.cpp}{70--86}. Node $c$ is \emph{eligible} for slot $\tslot$ iff
$\indic{H(c,\tslot) \le T_\tslot} = 1$, where $T_\tslot$ is the current \PoN difficulty target (below).
Every confirmed \FluxNode, of any tier, computes exactly one such hash per slot; eligibility does not
depend on tier or on collateral size beyond confirmation
\shipped \srccode{fluxd/\allowbreak{}src/\allowbreak{}pon/\allowbreak{}pon.cpp}{400--433}. This is sortition over the confirmed-node set, not
a lottery weighted by stake amount.
\end{definition}

The target $T_\tslot$ comes from \texttt{Get\allowbreak{}Next\allowbreak{}PONWork\allowbreak{}Required}. Before \PoN activation it is unused;
for the first \texttt{n\allowbreak{}Pon\allowbreak{}Difficulty\allowbreak{}Window} $= 30$ blocks (mainnet) after activation it is held fixed at \texttt{ponStartLimit},
which the source's own comment describes as giving each node roughly a 1-in-10{,}000 chance of
eligibility per slot \shipped \srccode{fluxd/\allowbreak{}src/\allowbreak{}chainparams.cpp}{104}. Thereafter every block
retargets: the actual timespan over the trailing 30-block window is clamped to
$\big[\tfrac{4}{5},\,\tfrac{5}{4}\big] \times (30-1)\spacingpon$ (i.e.\ $-20\%/{+}25\%$) and used,
Digishield-style, to scale the previous target
\shipped \srccode{fluxd/\allowbreak{}src/\allowbreak{}pon/\allowbreak{}pon.cpp}{149--157}. The protocol floor is \texttt{ponLimit}, described
in the same source comments as giving roughly a 1-in-16 minimum per-slot eligibility chance, mainnet
\shipped \srccode{fluxd/\allowbreak{}src/\allowbreak{}chainparams.cpp}{103}.

\begin{algorithm}[H]
\DontPrintSemicolon
\caption{\PoN minting loop, per node (one thread, wall-clock polled)}
\label{alg:pon-minter}
\KwIn{collateral outpoint $c$, \FluxNode private key $sk$, chain tip pointer $\mathrm{tip}$}
\KwOut{a signed candidate block submitted to the network, or no action this slot}
wait for network peers and end of initial block download\tcp*{\srccode{fluxd/\allowbreak{}src/\allowbreak{}pon/\allowbreak{}pon-minter.cpp}{121--136}}
\While{node is running}{
  \If{\PoN is active \emph{and} this node's own confirmation status is \emph{confirmed}}{
    $t \gets$ current wall-clock time\;
    $\tslot \gets \tslot(t)$ (Definition~\ref{def:pon-eligibility})\tcp*{\srccode{fluxd/\allowbreak{}src/\allowbreak{}pon/\allowbreak{}pon-minter.cpp}{167--177}}
    \If{$\tslot \neq$ \textup{lastAttemptedSlot} \emph{and} $t - \mathrm{tip.nTime} \ge 5\,\mathrm{s}$}{
      compute eligibility and rank of $c$ over all confirmed nodes for slot $\tslot$\tcp*{\srccode{fluxd/\allowbreak{}src/\allowbreak{}pon/\allowbreak{}pon-minter.cpp}{222--236}}
      \If{$c$ eligible for $\tslot$}{
        sleep $\mathrm{rank}(c) \times 4000\,\mathrm{ms}$\tcp*{orphan-avoidance heuristic only --- no consensus enforcement, \srccode{fluxd/\allowbreak{}src/\allowbreak{}pon/\allowbreak{}pon-minter.cpp}{78--104,217--221,250}}
        \If{$\tslot$ has not expired \emph{and} $\mathrm{tip}$ has not moved}{
          build candidate block, coinbase paying the dev-fund address\tcp*{\srccode{fluxd/\allowbreak{}src/\allowbreak{}pon/\allowbreak{}pon-minter.cpp}{288--314}}
          sign block hash with $sk$ via \textsc{SignPONBlock}\;
          submit via \textsc{ProcessNewBlock}\tcp*{\srccode{fluxd/\allowbreak{}src/\allowbreak{}pon/\allowbreak{}pon-minter.cpp}{29--76}}
        }
      }
      mark $\tslot$ as \textup{lastAttemptedSlot}\;
    }
  }
  sleep 1 s (poll interval)\tcp*{\srccode{fluxd/\allowbreak{}src/\allowbreak{}pon/\allowbreak{}pon-minter.cpp}{175}}
}
\end{algorithm}

\emph{Invariant.} A node attempts at most one block per slot: \texttt{last\allowbreak{}Attempted\allowbreak{}Slot} is checked
before eligibility is computed and set on every exit from the slot, so a repeated poll of the same slot is a no-op
\shipped \srccode{fluxd/\allowbreak{}src/\allowbreak{}pon/\allowbreak{}pon-minter.cpp}{167--177}.

\emph{Failure mode.} Only \texttt{std::\allowbreak{}runtime\_error} is caught at the top of the loop, and on
catching it the thread logs once and returns, with no restart logic
\shipped \srccode{fluxd/\allowbreak{}src/\allowbreak{}pon/\allowbreak{}pon-minter.cpp}{343--346}; any other exception
type escapes the bare \texttt{std::thread} body \srccode{fluxd/\allowbreak{}src/\allowbreak{}pon/\allowbreak{}pon-minter.cpp}{356}
and, under C++ semantics, calls \texttt{std::terminate}, aborting the whole daemon rather than only
the minting thread.

The rank-based delay in the algorithm's inner branch is worth stating plainly, because "Proof of Node"
invites a reading in which rank determines who is allowed to produce a block. It does not. The
$\mathrm{rank}(c) \times 4000\,\mathrm{ms}$ wait is documented in the source itself as an
orphan-reduction heuristic: a lower-ranked (slower) eligible node waits longer before broadcasting, so
that a higher-ranked node's block has time to propagate and preempt it. No peer enforces this delay;
any validly-signed block for an eligible node's slot is accepted on ordinary first-seen/most-work
rules regardless of the sender's rank
\shipped \srccode{fluxd/\allowbreak{}src/\allowbreak{}pon/\allowbreak{}pon-minter.cpp}{78--104,217--221}.

Consensus does not pin the slot a block claims to wall-clock time either: the slot is derived from
the block's own \texttt{nTime} \shipped \srccode{fluxd/\allowbreak{}src/\allowbreak{}pon/\allowbreak{}pon.cpp}{63--67}, and
the only rules on \texttt{nTime} are that it exceed the parent's
\srccode{fluxd/\allowbreak{}src/\allowbreak{}main.cpp}{5564--5567} and lie at most 300\,s ahead of the validator's
adjusted time \srccode{fluxd/\allowbreak{}src/\allowbreak{}main.cpp}{5368--5371}. A producer may therefore evaluate
$H(c,\tslot)$ over every slot in that window --- about eleven at 30\,s spacing --- and publish at once
for any slot in which one of its nodes qualifies, whereas the minting loop of
Algorithm~\ref{alg:pon-minter} tries only the current slot.

\subsection{Fork choice}
\label{sec:consensus-forkchoice}

\begin{property}[\PoN chain-work is a block count]
\label{prop:pon-forkchoice}
For every block with $\mathrm{nVersion} \ge 100$ (every \PoN block), \texttt{GetBlockProof} returns the
fixed value $2^{40}$, independent of the block's actual difficulty target
\shipped \srccode{fluxd/\allowbreak{}src/\allowbreak{}pow.cpp}{284--290}. Consequently, for two competing \PoN chains of length
$n_1$ and $n_2$ blocks since their most recent common ancestor, cumulative chain work is exactly
$n_1 \cdot 2^{40}$ versus $n_2 \cdot 2^{40}$: "most cumulative work" fork choice is equivalent to "most
blocks," and among chains of equal length, the deciding factor is which block a node saw and connected
first, not any comparison of cryptographic hardness
\shipped \srccode{fluxd/\allowbreak{}src/\allowbreak{}pow.cpp}{284--290}.
\end{property}

\begin{proof}[Proof sketch]
Chain work under the inherited Bitcoin/Zcash chain-selection code is the sum of \texttt{GetBlockProof}
over every block on a candidate chain since the fork point. \texttt{GetBlockProof} is defined
piecewise on \texttt{nVersion}; for $\mathrm{nVersion} \ge 100$ it returns the constant
\texttt{arith\_uint256(1)} $\ll 40$ regardless of \texttt{nBits}
\shipped \srccode{fluxd/\allowbreak{}src/\allowbreak{}pow.cpp}{284--290}. Summing a constant $n$ times over a chain of $n$
\PoN blocks gives $n \cdot 2^{40}$ exactly, with no dependence on the per-block difficulty target
computed in \S\ref{sec:consensus-sortition}. Comparing $n_1 \cdot 2^{40}$ against $n_2 \cdot 2^{40}$ is
therefore equivalent to comparing $n_1$ against $n_2$; the constant factor cancels. When $n_1 = n_2$,
the comparison is a tie under this metric, and the surviving chain is whichever block the node's own
chain-selection logic connected first, an ordering determined by network propagation and arrival time
rather than by proof of work.
\end{proof}

The consequence is not a defect to be patched quietly; it is a different security argument than the
one "proof of work" trains a reader to expect. \PoN's resistance to a chain reorganization does not
come from the cost of recomputing accumulated hashing work. It comes from the cost of controlling
enough collateral to win sortition (Definition~\ref{def:pon-eligibility}) across enough consecutive
slots to outpace and outlast the honest chain, and from how quickly an alternate chain of blocks can be
produced and propagated — a length race, not a work race. A reader who assumes Nakamoto-style
work-accumulation will draw the wrong conclusion about what a reorg of a given depth costs an attacker
on \PoN.

\begin{remark}
\texttt{n\allowbreak{}Minimum\allowbreak{}Chain\allowbreak{}Work}, the anti-eclipse floor consulted during initial block download, was last
updated to match the chain work associated with \texttt{UPGRADE\_KAMIOOKA}
\shipped \srccode{fluxd/\allowbreak{}src/\allowbreak{}chainparams.cpp}{168--170}; it has not been revised for the fixed-$2^{40}$
-per-block regime that \PoN introduced. Whether a floor computed under the old work-accumulation model
still functions as intended against a fake alternate chain once every block after height
2{,}020{,}000 contributes the same fixed increment is not established by the source \srcinferred.
\end{remark}

\subsection{The \FluxNode lifecycle}
\label{sec:consensus-lifecycle}

A deterministic \FluxNode moves through the states
\texttt{FLUXNODE\_TX\_ERROR}(0) $\to$ \texttt{FLUXNODE\_TX\_STARTED}(1) $\to$
\texttt{FLUXNODE\_TX\_DOS\_PROTECTION}(2) / \texttt{FLUXNODE\_TX\_CONFIRMED}(3) $\to$
\texttt{FLUXNODE\_TX\_MISS\_CONFIRMED}(4) / \texttt{FLUXNODE\_TX\_EXPIRED}(5)
\shipped \srccode{fluxd/\allowbreak{}src/\allowbreak{}fluxnode/\allowbreak{}fluxnode.h}{109--116}.

\begin{itemize}
\item \textbf{Entry (START).} The operator broadcasts a \texttt{FLUXNODE\_START\_TX\_TYPE} transaction
spending a collateral UTXO of a valid tier amount, signed by the collateral key
\shipped \srccode{fluxd/\allowbreak{}src/\allowbreak{}fluxnode/\allowbreak{}activefluxnode.cpp}{253--288}. Consensus requires the collateral
to have $\ge 100$ confirmations (\texttt{FLUXNODE\_MIN\_CONFIRMATION\_DETERMINISTIC})
\shipped \srccode{fluxd/\allowbreak{}src/\allowbreak{}fluxnode/\allowbreak{}fluxnode.h}{20} and to match a tier's valid amount
\shipped \srccode{fluxd/\allowbreak{}src/\allowbreak{}coins.cpp}{630--686}. The node enters \texttt{map\allowbreak{}Start\allowbreak{}Tx\allowbreak{}Tracker}
\shipped \srccode{fluxd/\allowbreak{}src/\allowbreak{}fluxnode/\allowbreak{}fluxnode.cpp}{304--352}.

\item \textbf{STARTED $\to$ DOS\_PROTECTION.} If no confirm transaction arrives within
\texttt{FLUXNODE\_START\_TX\_EXPIRATION\_HEIGHT} $= 60$ blocks (240 post-\PoN,
\texttt{\_V2}) of the start height, the node moves to \texttt{map\allowbreak{}Start\allowbreak{}Tx\allowbreak{}DOSTracker}
\shipped \srccode{fluxd/\allowbreak{}src/\allowbreak{}fluxnode/\allowbreak{}fluxnode.cpp}{479--508}; it is fully purged once
\texttt{FLUXNODE\_DOS\_REMOVE\_AMOUNT} $= 180$ blocks (720 post-\PoN, \texttt{\_V2}) have elapsed since
the \emph{start} height, i.e.\ 120 (480 post-\PoN) blocks after entering the DoS list
\shipped \srccode{fluxd/\allowbreak{}src/\allowbreak{}fluxnode/\allowbreak{}fluxnode.cpp}{205--229},
\srccode{fluxd/\allowbreak{}src/\allowbreak{}fluxnode/\allowbreak{}fluxnode.h}{26--28}.

\item \textbf{STARTED $\to$ CONFIRMED.} The operator builds a \texttt{FLUXNODE\_CONFIRM\_TX\_TYPE}
(\texttt{INITIAL\_CONFIRM}), which is benchmark-signed by \fluxbench before wallet broadcast (see
\S\ref{sec:consensus-benchmark}) \shipped \srccode{fluxd/\allowbreak{}src/\allowbreak{}fluxnode/\allowbreak{}activefluxnode.cpp}{380--407,18--74}.
On chain acceptance the node enters \texttt{map\allowbreak{}Confirmed\allowbreak{}Fluxnode\allowbreak{}Data} and is inserted at the
\emph{front} of its tier's payment list, which is then re-sorted
\shipped \srccode{fluxd/\allowbreak{}src/\allowbreak{}fluxnode/\allowbreak{}fluxnode.cpp}{1502--1511,1466--1474}.

\item \textbf{Re-confirmation (CONFIRMED, periodic).} \texttt{Check\allowbreak{}If\allowbreak{}Needs\allowbreak{}Next\allowbreak{}Confirm} triggers a
\texttt{UPDATE\_CONFIRM} once $\hgt - \texttt{n\allowbreak{}Last\allowbreak{}Confirmed\allowbreak{}Block\allowbreak{}Height}$ exceeds
\texttt{FLUXNODE\_CONFIRM\_UPDATE\_MIN\_HEIGHT\_V1} $= 40$ / \texttt{\_V2} $= 120$ /
\texttt{\_V3} $= 500$ (\PoN), by active upgrade
\shipped \srccode{fluxd/\allowbreak{}src/\allowbreak{}fluxnode/\allowbreak{}fluxnode.cpp}{669--688}.

\item \textbf{CONFIRMED $\to$ EXPIRED.} If no re-confirmation lands within
\texttt{FLUXNODE\_CONFIRM\_UPDATE\_EXPIRATION\_HEIGHT\_V1} $= 60$ / \texttt{\_V2} $= 80$ /
\texttt{\_V3} $= 160$ / \texttt{\_V4} $= 640$ (\PoN) blocks of \texttt{n\allowbreak{}Last\allowbreak{}Confirmed\allowbreak{}Block\allowbreak{}Height}, the
node is removed from the paid set \shipped \srccode{fluxd/\allowbreak{}src/\allowbreak{}fluxnode/\allowbreak{}fluxnode.cpp}{1716--1732},
\srccode{fluxd/\allowbreak{}src/\allowbreak{}fluxnode/\allowbreak{}fluxnode.h}{31--34}.

\item \textbf{CONFIRMED $\to$ removed (collateral spent).} If the confirmation collateral UTXO is
spent, the node is removed \shipped \srccode{fluxd/\allowbreak{}src/\allowbreak{}fluxnode/\allowbreak{}fluxnode.cpp}{285--290}.

\item \textbf{CONFIRMED $\to$ removed (collateral migration mismatch).} During the V1$\to$V2
collateral-amount transition window (\S\ref{sec:consensus-tiers}), a confirmed node whose collateral no
longer matches the \emph{new} valid amount for its tier is removed once that tier's transition-end
height is reached \shipped \srccode{fluxd/\allowbreak{}src/\allowbreak{}fluxnode/\allowbreak{}fluxnode.cpp}{231--283}.
\end{itemize}

\textbf{Payment queue ordering.} \texttt{Fluxnode\allowbreak{}List\allowbreak{}Data::\allowbreak{}operator<} orders each tier's list by
\texttt{n\allowbreak{}Last\allowbreak{}Paid\allowbreak{}Height} (or \texttt{n\allowbreak{}Confirmed\allowbreak{}Block\allowbreak{}Height} if never paid) ascending, tie-broken by
whether \texttt{n\allowbreak{}Last\allowbreak{}Paid\allowbreak{}Height} is set and then by outpoint
\shipped \srccode{fluxd/\allowbreak{}src/\allowbreak{}fluxnode/\allowbreak{}fluxnode.h}{313--328} — a strict longest-time-since-last-paid
FIFO per tier. \texttt{GetNextPayment} reads the front of the tier's list, skipping and popping any
stale (unconfirmed) entries it finds there
\shipped \srccode{fluxd/\allowbreak{}src/\allowbreak{}fluxnode/\allowbreak{}fluxnode.cpp}{713--778}.

\begin{figure}[H]
\centering
\begin{tikzpicture}[
  font=\scriptsize,
  st/.style={draw, rounded corners=2pt, align=center, minimum height=8mm, text width=26mm,
             inner xsep=2pt},
  term/.style={st, fill=black!8},
  live/.style={st, very thick, fill=black!4},
  arr/.style={-{Latex[length=1.8mm]}, thick},
  lbl/.style={font=\tiny, align=center, inner sep=1.5pt},
]
  \node[st]                                (start) {\texttt{TX\_STARTED}};
  \node[st,   below=24mm of start]         (dos)   {\texttt{TX\_DOS\_}\\\texttt{PROTECTION}};
  \node[live, right=40mm of start]         (conf)  {\texttt{TX\_CONFIRMED}};
  \node[term, right=40mm of dos]           (exp)   {\texttt{TX\_EXPIRED}};
  \node[st,   above=18mm of conf]          (miss)  {\texttt{TX\_MISS\_}\\\texttt{CONFIRMED}};
  \node[term, left=28mm of start]          (err)   {\texttt{TX\_ERROR}};

  \draw[arr] (err) -- node[lbl, above] {start tx,\\collateral aged\\$\ge 100$ blocks} (start);
  \draw[arr] (start) -- node[lbl, left, pos=0.5] {no confirm within\\60 / 240 (\PoN)\\blocks} (dos);
  \draw[arr] (start) -- node[lbl, above] {confirm tx\\(\textsc{initial})} (conf);
  \draw[arr] (dos) -- node[lbl, below] {purged 180 / 720 (\PoN) blocks\\after the start height} (exp);

  %% Self-loop kept on the right of TX_CONFIRMED so it clears TX_MISS_CONFIRMED above.
  \draw[arr] (conf) to[out=-25, in=25, looseness=7]
      node[lbl, right, pos=0.5, text width=27mm] {re-confirm every\\40 / 120 / 500 (\PoN) blocks} (conf);
  %% Both vertical edge labels sit on the left so nothing collides with the loop.
  \draw[arr] (conf) -- node[lbl, left, pos=0.5, text width=30mm]
      {no re-confirm within\\60 / 80 / 160 / 640 (\PoN) blocks;\\or collateral spent;\\or V1$\to$V2 amount mismatch} (exp);
  \draw[arr] (conf) -- node[lbl, left, pos=0.5] {missed} (miss);
  \draw[arr, dashed] (miss) to[out=180, in=90] node[lbl, above, pos=0.35] {re-confirm} (start);

  %% Spans the full width below both bottom states rather than sitting under TX_CONFIRMED,
  %% where it used to print on top of the conf-to-exp edge label.
  \node[lbl, below=13mm of $(dos)!0.5!(exp)$, text width=112mm]
    {on entry to \texttt{TX\_CONFIRMED}: inserted at the \emph{front} of its tier's payment list, then the
     list is re-sorted; \texttt{GetNextPayment} pops the front of the tier-scoped FIFO ordered by
     longest-time-since-paid};
  \end{tikzpicture}
\caption{The \FluxNode lifecycle state machine: states, transitions, and their governing height
constants \shipped \srccode{fluxd/\allowbreak{}src/\allowbreak{}fluxnode/\allowbreak{}fluxnode.h}{109--116}. Where two or more values are
shown, they are the successive versions of the constant selected by the active network upgrade, the
last being the \PoN-era value \srccode{fluxd/\allowbreak{}src/\allowbreak{}fluxnode/\allowbreak{}fluxnode.h}{23--44}. Only
\texttt{TX\_CONFIRMED} is in the paid set.}
\label{fig:fluxnode-lifecycle}
\end{figure}

\subsection{Tiers and collateral}
\label{sec:consensus-tiers}

\begin{table}[H]
\centering
\caption{\FluxNode collateral by tier and version, mainnet. Confirmation requirement is uniform across
tiers and versions.}
\label{tab:tiers}
\begin{tabular}{lrrll}
\toprule
Tier & V1 collateral & V2 collateral & Transition window $[\hgt_{\mathrm{start}},\hgt_{\mathrm{end}})$ & Confirmation \\
\midrule
Cumulus ($\tC$) & $\coll{\tC}=10{,}000\flux$ & $1{,}000\flux$ & $[1{,}076{,}532,\ 1{,}086{,}612)$ & 100 blocks \\
Nimbus ($\tN$)  & $\coll{\tN}=25{,}000\flux$ & $12{,}500\flux$ & $[1{,}081{,}572,\ 1{,}092{,}372)$ & 100 blocks \\
Stratus ($\tS$) & $\coll{\tS}=100{,}000\flux$ & $40{,}000\flux$ & $[1{,}087{,}332,\ 1{,}097{,}412)$ & 100 blocks \\
\bottomrule
\end{tabular}
\end{table}
\shipped \srccode{fluxd/\allowbreak{}src/\allowbreak{}fluxnode/\allowbreak{}fluxnode.h}{57--63} (V1/V2 amounts),
\srccode{fluxd/\allowbreak{}src/\allowbreak{}chainparams.cpp}{308--315} (transition windows),
\srccode{fluxd/\allowbreak{}src/\allowbreak{}fluxnode/\allowbreak{}fluxnode.h}{20} (100-block confirmation requirement,
\texttt{FLUXNODE\_MIN\_CONFIRMATION\_DETERMINISTIC}). Within a transition window, both the old and new
collateral amounts for that tier are accepted \shipped \srccode{fluxd/\allowbreak{}src/\allowbreak{}coins.cpp}{774--809}. The V2
amounts match the collateral recorded for every one of the 6{,}489 currently registered \FluxNodes in
the public registry \srcmeasured{api.runonflux.io/\allowbreak{}daemon/\allowbreak{}getzelnodecount}{2026-09-03}, i.e.\ the
migration is complete on mainnet as measured.

\subsection{Benchmarking and eligibility}
\label{sec:consensus-benchmark}

Confirming or re-confirming a \FluxNode requires a benchmark-signed transaction produced by
\fluxbench, a companion daemon that runs on the same host as \fluxd. \fluxbench measures CPU
throughput with \texttt{sysbench cpu -\allowbreak{}-\allowbreak{}cpu-\allowbreak{}max-\allowbreak{}prime=60000}
\shipped \srccode{fluxbench/\allowbreak{}src/\allowbreak{}fluxnode/\allowbreak{}benchmarks.cpp}{2066--2364}; disk write speed with three
\texttt{dd} runs per target, dropping the worst and averaging the remaining two, floored at
40 MB/s for reporting \shipped \srccode{fluxbench/\allowbreak{}src/\allowbreak{}fluxnode/\allowbreak{}benchmarks.h}{602--615,1020--1104},
\srccode{fluxbench/\allowbreak{}src/\allowbreak{}fluxnode/\allowbreak{}benchmarks.cpp}{1936}; bandwidth with the Ookla \texttt{speedtest} CLI,
combined with an \texttt{iftop} sample of bandwidth already in use
\shipped \srccode{fluxbench/\allowbreak{}src/\allowbreak{}fluxnode/\allowbreak{}benchmarks.cpp}{1358--1465,1980--2032}; and RAM headroom with
a \texttt{stress-ng} allocation test sized 8G/4G/2G for Stratus/Nimbus/Cumulus
\shipped \srccode{fluxbench/\allowbreak{}src/\allowbreak{}fluxnode/\allowbreak{}benchmarks.cpp}{3649--3667}.

\begin{table}[H]
\centering
\caption{Per-tier hardware thresholds enforced by \texttt{Validate\allowbreak{}Benchmark\allowbreak{}Requirements}, mainnet.}
\label{tab:benchmark-thresholds}
\begin{tabular}{lrrrrrrr}
\toprule
Tier & Cores & Threads & RAM & Storage & EPS & Disk write & Up/Down \\
\midrule
Cumulus & 2 & 4 & 7 GB & 220 GB & 240 & 180 MB/s & 25/25 Mbit/s \\
Nimbus  & 4 & 8 & 30 GB & 440 GB & 640 & 180 MB/s & 50/50 Mbit/s \\
Stratus & 8 & 16 & 61 GB & 880 GB & 1{,}520 & 400 MB/s & 100/100 Mbit/s \\
\bottomrule
\end{tabular}
\end{table}
\shipped \srccode{fluxbench/\allowbreak{}src/\allowbreak{}fluxnode/\allowbreak{}benchmarks.h}{58--84}. Cumulus's RAM requirement is 7 GB on
standard hardware and 3 GB on Jetson boards specifically
\shipped \srccode{fluxbench/\allowbreak{}src/\allowbreak{}fluxnode/\allowbreak{}benchmarks.h}{58--66}; measured bandwidth additionally
receives an additive tolerance bonus before the threshold check of $+2$/$+4$/$+9$ Mbit/s for
Cumulus/Nimbus/Stratus \shipped \srccode{fluxbench/\allowbreak{}src/\allowbreak{}fluxnode/\allowbreak{}benchmarks.h}{86--88}.

The result is signed by \texttt{BenchmarkSign}, which hashes the transaction's \texttt{sig},
\texttt{benchmarkTier}, \texttt{benchmark\allowbreak{}Sig\allowbreak{}Time}, and \texttt{ip} fields and signs with the fluxnode
key selected for the current time window
\shipped \srccode{fluxbench/\allowbreak{}src/\allowbreak{}fluxnode/\allowbreak{}benchmarks.cpp}{3439--3480}. The same verification routine,
\texttt{Check\allowbreak{}Benchmark\allowbreak{}Signature}, is linked into \fluxd and run as part of consensus validation of the
registration transaction \shipped \srccode{fluxbench/\allowbreak{}src/\allowbreak{}fluxnode/\allowbreak{}benchmarks.cpp}{3511--3522}. Three
further mechanisms can substitute a stored or historical value for a fresh measurement under specified
conditions — an enterprise-app short-circuit, a FluxOS-reported fallback, and a
single-thread-extrapolation fallback for the CPU test — and all three converge on the same
\texttt{Validate\allowbreak{}Benchmark\allowbreak{}Requirements()} acceptance gate
\shipped \srccode{fluxbench/\allowbreak{}src/\allowbreak{}fluxnode/\allowbreak{}benchmarks.cpp}{2178--2234,771--783,2883--3222}.

\paragraph{The property that does not hold.} On the \fluxd side, this same signature is verified
against \texttt{vec\allowbreak{}Benchmarking\allowbreak{}Public\allowbreak{}Keys} — five keys, each with a UNIX-timestamp validity window,
i.e.\ a rotating allowlist \shipped \srccode{fluxd/\allowbreak{}src/\allowbreak{}chainparams.cpp}{233--240}. These keys are
shared across the entire fleet: they identify "signed by something holding one of the five
network-wide benchmarking keys," not "signed by this specific operator's own key." \fluxd obtains the
benchmark result itself by shelling out — \texttt{popen()}-ing the local \texttt{fluxbench-cli}
process — and treats a returned \texttt{status:\allowbreak{}"complete"} field as sufficient to sign and broadcast
the registration transaction \shipped \srccode{fluxd/\allowbreak{}src/\allowbreak{}fluxnode/\allowbreak{}benchmarks.cpp}{74--102,227--291}.
The registration transaction must still match a valid on-chain collateral amount for the claimed tier
\shipped \srccode{fluxd/\allowbreak{}src/\allowbreak{}coins.cpp}{630--686}.

Put precisely: collateral binds the tier; the benchmark signature does not bind the hardware to the
tier it claims. An operator cannot register a tier they have not collateralized, but nothing beyond a
local, same-host process attests that the collateralized hardware in fact clears that tier's CPU/RAM
/disk/bandwidth thresholds in Table~\ref{tab:benchmark-thresholds}. \textbf{The correct claim is
benchmark-gating bypass, not tier forgery.} \S14 takes up the structurally identical problem one layer
up, for \ArcaneOS/\SAS attestation, and presents the planned bonded attestation network — where a bond
is put at stake rather than a shared secret being checked — as the mechanism intended to fix exactly
this property \planned \srcroadmap{timeline: new benchmark tool, signed node attestation}.

\subsection{The emergency-block path}
\label{sec:consensus-emergency}

A second, parallel block-authorization path exists outside \PoN sortition entirely, intended for
chain-halt recovery. A block whose \texttt{nodes\allowbreak{}Collateral} outpoint hash equals a fixed marker value
and whose index is 0 is treated as an \emph{emergency block}
\shipped \srccode{fluxd/\allowbreak{}src/\allowbreak{}emergencyblock.cpp}{24--30}. Such a block requires at least
\texttt{n\allowbreak{}Emergency\allowbreak{}Min\allowbreak{}Signatures} $= 2$ valid ECDSA signatures over \texttt{block.\allowbreak{}Get\allowbreak{}Hash()}, drawn from
a fixed set of 4 hardcoded public keys — a 2-of-4 multisignature, on mainnet and testnet alike
\shipped \srccode{fluxd/\allowbreak{}src/\allowbreak{}chainparams.cpp}{242--253}. Accepting an emergency block bypasses
\texttt{Check\allowbreak{}Proof\allowbreak{}Of\allowbreak{}Node} entirely \shipped \srccode{fluxd/\allowbreak{}src/\allowbreak{}pon/\allowbreak{}pon.cpp}{205--207} and bypasses the
per-node signature verification that an ordinary \PoN block requires
\shipped \srccode{fluxd/\allowbreak{}src/\allowbreak{}pon/\allowbreak{}pon.cpp}{236--252}.

The only gate on \emph{when} an emergency block may be produced is
\texttt{Is\allowbreak{}Emergency\allowbreak{}Block\allowbreak{}Allowed(n\allowbreak{}Height, n\allowbreak{}Time)}, which checks nothing beyond
\texttt{Is\allowbreak{}PONActive(n\allowbreak{}Height)} — there is no time-since-last-block check and no stall-detection logic
\shipped \srccode{fluxd/\allowbreak{}src/\allowbreak{}emergencyblock.cpp}{183--191} — despite a doc-comment on the type stating
"Can add additional restrictions like time delays"
\shipped \srccode{fluxd/\allowbreak{}src/\allowbreak{}emergencyblock.h}{41--44}. Concretely: any two holders of the four
emergency keys can mint a fully valid block at any height once \PoN is active, at any time, with no
on-chain evidence that the network was in fact stalled. This is a consensus-level centralization
surface, and it is stated here as one, not as a caveat.

The team intends to raise this threshold to \textbf{3-of-4 with a newly generated key set},
aligning consensus emergency-key governance with the bridge's 3-of-4 quorum (\S22), with no
activation height scheduled \planned \srcintent{design intake}. This is a stated direction of travel, not a deployed change: \texttt{n\allowbreak{}Emergency\allowbreak{}Min\allowbreak{}Signatures}
remains \textbf{2} against the current four-key set on mainnet and testnet today, which is why this
section states 2-of-4 \shipped \srccode{fluxd/\allowbreak{}src/\allowbreak{}chainparams.cpp}{242--253} above. Raising the
threshold does not change the shape of the disclosure just made: a 3-of-4 key set, like the 2-of-4
set it would replace, produces blocks outside \PoN{} eligibility and outside the per-node signature
check, with no time or stall-detection gating either before or after the change; what moves is only
the number of colluding or compromised keys required, from two of four to three of four. The absence
of any temporal or liveness constraint on \emph{when} an emergency block may be produced --- the more
consequential fact of the two --- is untouched by this parameter.

\subsection{What activation actually cost}
\label{sec:consensus-cost}

Activating \PoN required two emergency interventions visible directly in the code.
\texttt{MAX\_REORG\_LENGTH}, normally capped at 40 blocks
\shipped \srccode{fluxd/\allowbreak{}src/\allowbreak{}main.h}{74}, was overridden to 5{,}000 for heights
$[2{,}020{,}000,\ 2{,}025{,}000]$, commented "Temporary deep reorg limit during PON stabilization
period" \shipped \srccode{fluxd/\allowbreak{}src/\allowbreak{}main.cpp}{8983--8988}. A checkpoint cluster at heights
2{,}020{,}500; 2{,}021{,}000; 2{,}021{,}500; 2{,}022{,}000; and 2{,}029{,}000 was added, labelled "PON
activation and stabilization checkpoints - EMERGENCY FORK RECOVERY"
\shipped \srccode{fluxd/\allowbreak{}src/\allowbreak{}chainparams.cpp}{277--283}. \S\ref{sec:migration} returns to this as evidence for what a
consensus-level change actually requires of this network.

\subsection{Measured rotation}
\label{sec:consensus-rotation}

Under the front-of-list, longest-time-since-paid ordering described in
\S\ref{sec:consensus-lifecycle}, and given that each block pays one node from every tier, the expected
number of blocks between successive payments to a node in tier $x$ is $\rot{x} \defeq \ncount{x}$, the
tier's live node count. Table~\ref{tab:pon-rotation} compares this expectation against the maximum
observed gap in the public registry, chain tip height 2{,}917{,}115
\srcmeasured{api.runonflux.io/\allowbreak{}daemon/\allowbreak{}listzelnodes}{2026-09-03}.

\begin{table}[H]
\centering
\caption{Measured \PoN payment rotation vs. the deterministic per-tier model.}
\label{tab:pon-rotation}
\begin{tabular}{lrrrr}
\toprule
Tier & Nodes ($\ncount{\cdot}$) & Expected rotation ($\rot{\cdot}$, blocks) & Max blocks since paid & Duration at $\spacingpon$ \\
\midrule
Cumulus & 3{,}247 & 3{,}247 & 3{,}202 & 27.06 h \\
Nimbus  & 1{,}615 & 1{,}615 & 1{,}597 & 13.46 h \\
Stratus & 1{,}627 & 1{,}627 & 1{,}631 & 13.56 h \\
\bottomrule
\end{tabular}
\end{table}
\srcmeasured{api.runonflux.io/\allowbreak{}daemon/\allowbreak{}listzelnodes}{2026-09-03}. Stratus's observed maximum exceeds its
own node count by 4 blocks; this is consistent with the stale-entry skip loop in
\texttt{GetNextPayment} (\S\ref{sec:consensus-lifecycle}) rather than a violation of the FIFO ordering
rule. This measured rotation is the empirical confirmation on which \S7's fairness proposition for
\PNR depends: that proposition rests on the determinism of the \PoN \emph{payment} rotation, not on
placement determinism, and the two must not be conflated.

%% Drain this section's float queue before the next begins.
\FloatBarrier
